\documentclass[12pt]{article}

\usepackage[margin=1in]{geometry}
\usepackage[T1]{fontenc}
\usepackage[utf8]{inputenc}
% microtype omitted (requires scalable fonts not available in this installation)
\usepackage{amsmath, amssymb, mathtools}
\usepackage{graphicx}
\usepackage{booktabs}
\usepackage{natbib}
\usepackage{setspace}
\usepackage[hidelinks]{hyperref}
\newcommand{\orcid}[1]{\href{https://orcid.org/#1}{ORCID: #1}}
\usepackage{csquotes}
\usepackage{siunitx}
\usepackage{subcaption}
\usepackage{tikz}
\usetikzlibrary{arrows.meta,positioning,fit,calc,shadows.blur}
\usepackage{tabularx}
\newcommand{\tablenote}[1]{%
\vspace{0.35em}
\begin{flushleft}
\footnotesize\emph{Note}. #1
\end{flushleft}
}

\usepackage{fancyhdr}
\pagestyle{fancy}
\fancyhf{}
\rhead{\thepage}
\renewcommand{\headrulewidth}{0pt}
\setlength{\headheight}{15pt}

\doublespacing
\graphicspath{{figures/}{figures/}{./}}

\title{Angular profiling of contextual structure in implicit-association latency data}

\author{%
Felipe de Oliveira Matos\textsuperscript{1}
\quad
Marlos Andrade de Lima\textsuperscript{2}
\\[0.25em]
Cristian Zanon\textsuperscript{2}
\quad
Antonio Pereira\textsuperscript{3,*}
\\[0.9em]
\footnotesize
\begin{tabular}{@{}p{0.88\textwidth}@{}}
\centering
\textsuperscript{1}Department of Physiological Sciences, Universidade Estadual de Maringá,
87020-900, Maringá, PR, Brazil. \\[0.35em]
\centering
\textsuperscript{2}Department of Developmental Psychology and Personality,
Universidade Federal do Rio Grande do Sul, Porto Alegre, RS, Brazil. \\[0.35em]
\centering
\textsuperscript{3}Signal Processing Laboratory, Universidade Federal do Pará (UFPA),
66075-110, Belém, PA, Brazil. \\[0.75em]
\centering
Felipe de Oliveira Matos:
\href{mailto:fomatos@uem.br}{fomatos@uem.br};
\orcid{0000-0002-4926-4694}. \\[0.35em]
\centering
Marlos Andrade de Lima:
\href{mailto:andradelimarlos@gmail.com}{andradelimarlos@gmail.com};
\orcid{0000-0002-8901-2272}. \\[0.35em]
\centering
Cristian Zanon:
\href{mailto:crstn.zan@gmail.com}{crstn.zan@gmail.com};
\orcid{0000-0003-3822-5275}. \\[0.35em]
\centering
\textsuperscript{*}Corresponding author: Antonio Pereira,
\href{mailto:apereira@ufpa.br}{apereira@ufpa.br};
\orcid{0000-0002-0808-1058}.
\end{tabular}%
}

\date{}

\begin{document}
\maketitle

\begin{abstract}
Large online behavioral studies generate trial-level latency records at a scale that exceeds conventional scalar summaries. We introduce an angular-profiling workflow for describing task-level latency geometry in Implicit Association Test (IAT) data. The method aggregates trial latencies into participant-level binned trajectories, fits a three-term rotational basis with closed-form participant-level coefficients, and profiles a task-level angle, \(\theta\), as a compact descriptor of contextual separation. Smaller angles index greater contextual overlap in aggregated latency geometry, and larger angles index greater contextual separation.

Across public Project Implicit Gender--Science, Sexuality, and Age IATs, the profiled angle occupied a narrow low-angle regime and remained stable across binning choices, participant bootstrap resampling, and D-score-like strata. An independent life-satisfaction IAT produced a low-angle modal estimate with broader uncertainty, matching its smaller sample and different implementation. Predictive and calibration diagnostics defined the estimator's operational scope: smooth additive baselines matched local RMSE, while residual and distributional diagnostics favored the angular basis, and simulation-based calibration showed limited low-angle recoverability under the current sparse-binning estimator. The workflow provides a reproducible diagnostic layer for large-scale IAT comparison, quality control, and exploratory analysis of latency structure.
\end{abstract}

\paragraph{Keywords:} Implicit Association Test, response latency, profile likelihood, contextual separation, estimator calibration, reproducible behavioral methods

\section{Introduction}

Large-scale online behavioral studies increasingly produce trial-level latency records from hundreds of thousands of participants. These datasets create an opportunity for population-level measurement that goes beyond single effect-size summaries, but they also pose a methodological problem: how should dynamic latency structure be compressed without treating a convenient summary as an individual-level latent trait? The Implicit Association Test (IAT) is a useful test case because it produces robust latency asymmetries between congruent and incongruent pairings, while also generating repeated trial-level observations that can be aggregated into participant-level time courses \citep{greenwald1998,greenwald2003,kurdi2020iat_response}.

Here, dynamic latency structure refers to the organized shape of response-time trajectories across trials, blocks, and pairing contexts after individual responses are aggregated into participant-level curves. This information is not available from scalar self-report ratings and is only partly represented by conventional IAT effect-size summaries. The target of the present framework is therefore task-level latency geometry, not individual-level attitude classification.

Standard IAT analyses, including D-score summaries, mixed-effects models, and diffusion formulations, typically decompose response latencies into additive contributions from congruency, learning, and noise \citep{klauer2007process,ratcliff2016diffusion}. These approaches are useful for effect-size estimation, process decomposition, and prediction. They are less directly suited to summarizing how strongly two task contexts separate or overlap in the aggregated latency geometry. This distinction matters because psychometric debates over the IAT have repeatedly shown that individual-level construct interpretation is fragile and depends on task structure, stimulus materials, response modality, and aggregation level \citep{schimmack2019iat,kurdi2020iat_response,hogenboom2024stimuli,baranan2024facial,axt2024predictive}.

We introduce a profiling framework for estimating a low-dimensional task-level summary of contextual separation in IAT latency data. The method is designed for a specific measurement problem: standard summaries such as the D score capture mean congruency differences, but they do not directly describe whether the aggregated latency trajectories generated by different task contexts occupy overlapping or separated geometric regimes. The proposed rotational basis summarizes binned participant-level latency curves with a single profiled angle, \(\theta\). In this framework, smaller angles correspond to stronger contextual overlap, whereas larger angles correspond to stronger contextual separation. The rotational parameterization provides a compact basis for departures from simple additivity and yields a tractable profiled estimator. We evaluate \(\theta\) as a task-level descriptor of aggregated behavioral latency structure by testing its stability, task sensitivity, predictive behavior, and simulation-based calibration.

The angular basis is related to geometric and quantum-inspired models of context dependence, where context changes can be represented as rotations between partially overlapping regimes \citep{busemeyer2012,pothos2013,pothos2022quantum,trueblood2011}. In the present application, this formalism is used as a measurement tool for aggregated latency structure. Its value is assessed empirically through stability, task sensitivity, predictive diagnostics, null comparisons, and simulation-based calibration.

We first apply this framework to the Gender--Science IAT as a large public reference corpus with standard Project Implicit task structure and extensive participant coverage. We then examine whether the same descriptive geometry generalizes to two additional conventional public social-category IATs: a matched Sexuality IAT and the Age IAT. These comparisons preserve the broad Project Implicit architecture while changing the social category under study, allowing us to test whether the low-angle pooled geometry is domain-specific or a more general feature of conventional public IAT tasks. Finally, we analyze an independent self-referential life-satisfaction IAT administered outside the Project Implicit platform. This dataset differs in platform, protocol, language, sample size, stimulus format, and construct target, and is therefore treated as an external generalization case rather than as a matched domain comparison.

The analysis evaluates the proposed profiling framework along six methodological dimensions: (1) whether a single profiled angle can summarize aggregated latency dynamics at the population level; (2) whether the estimate is stable under binning choices and participant resampling; (3) whether it remains distinct from conventional D-score-like strata; (4) whether it generalizes across public social-category IAT domains; (5) whether a smaller independent self-referential IAT remains compatible with the same low-angle regime; and (6) how the profiled structural objective relates to temporal hold-out prediction, additive baselines, null controls, and simulation-based calibration.

The contribution is a reusable analysis workflow for behavioral latency data. The workflow constructs participant-level binned trajectories, estimates a profiled angular descriptor, compares the descriptor across task domains, and evaluates the estimator with binning sensitivity, participant bootstrap resampling, predictive checks, additive nulls, temporal permutations, and simulation-based calibration. The profiled angle is a task-level descriptor of aggregated latency geometry under a specified basis.

The empirical question is whether \(\theta\) provides stable compression of task-level contextual separation after accounting for binning choices, D-score-like stratification, additive null structure, temporal order, and task-domain changes. The calibration question is whether the same estimator recovers known generating angles under controlled simulation. The analyses below show a stable and task-sensitive pooled angular profile, together with a clear low-angle calibration boundary under the present sparse-binning estimator.

The resulting contribution combines five elements: (i) participant-level binning of trial latencies, (ii) a closed-form inner least-squares fit, (iii) one-dimensional profiling of a task-level angular summary, (iv) cross-domain comparison across large public IAT datasets and an independent self-referential IAT, and (v) diagnostic checks linking stability, prediction, residual structure, distributional fit, null behavior, and simulation-based recovery.

\section{Methods}

\subsection{Data and preprocessing}
We merged trial-level logs (\texttt{iat.txt}, \texttt{iat2.txt}, and \texttt{iat.0003.txt}) from the publicly released Project Implicit demo archive, specifically the \texttt{scienceiat} task, and isolated the critical blocks (3, 4, 6, and 7). For each participant and block, we normalized the trial index to $x\in[0,1]$ and aggregated latencies into $N_{\text{bins}}\in\{4,6,8\}$ quantile bins (our main analysis relies on $6$ bins). We excluded participants with fewer than 3 populated bins, producing individual curves $\{(x_t,y_t)\}_{t=1}^{T_i}$ where $T_i\in[3,8]$.

\paragraph{Analysis unit.}
The primary modeling unit is a binned latency curve. In the main pooled analyses, each retained participant contributes a participant-level trajectory after trial-level latencies have been aggregated into quantile bins. In blockwise analyses, each retained participant--block combination contributes a block-specific trajectory. Participant counts and curve counts are therefore reported separately when they refer to different analytic objects. The pooled criterion sums residuals over curves, while the bootstrap resamples participants.

\paragraph{Leak-free standardization.}
In hold-out experiments, we standardize latencies (z-scoring) \emph{per participant using only the training bins}, applying the resulting transformation to both training and test data to prevent outcome leakage. When estimating the global posterior for $\theta$ (which does not involve predictive splitting), we z-score across all bins to align the data with our Gaussian observation model.

\subsection{Reproducible analysis workflow}

The analysis was implemented as a scripted, reproducible pipeline. The pipeline ingests trial-level task files, standardizes domain-specific columns, applies latency cleaning and critical-block selection, constructs participant-level binned curves, profiles \(\theta\) over a fixed grid, and exports diagnostic tables, figures, and aggregate cross-domain summaries. Public Project Implicit datasets were processed from trial-level logs, and the Age IAT archive was downloaded automatically from OSF before preprocessing. The life-satisfaction IAT was processed from the original Testable exports. Manuscript-facing outputs were generated directly from the processed analysis files.

\subsection{Pooled angular-profile model}
We model standardized mean latency as
\begin{equation}
y_i(x)\;=\;a_i + b_i\cos(\omega x) + c_i\sin(\omega x) + \varepsilon,
\qquad \varepsilon\sim\mathcal{N}(0,\sigma^2),\quad
\omega=\tfrac{\pi}{180}\,\theta.
\end{equation}
For any fixed $\theta$, the coefficients $(a_i,b_i,c_i)$ have the closed-form OLS solution $\hat\beta_i(\theta)=(X_i^\top X_i)^{-1}X_i^\top y_i$ with design columns $1,\cos(\omega x),\sin(\omega x)$.

\subsection{Identifiability, pooling, and inference}

Because the functions $\{1,\cos(\omega x),\sin(\omega x)\}$ are linearly independent for $\omega \neq 0$, the participant-level design matrix has full rank whenever at least three distinct bins are observed. This formal statement, however, is weaker than practical identifiability. In the empirically relevant low-angle regime, the small-angle expansion given below implies near-collinearity between the intercept and cosine columns and weak effective curvature in the profiled objective. The resulting estimator can therefore be sharply concentrated at the pooled level while still exhibiting poor latent recovery under realistic noise.

We do not assume that a single common angle is a literal generative truth for all participants or all task blocks. Instead, the pooled estimator
\[
\hat{\theta}_{\mathrm{pool}}=\arg\min_\theta RSS(\theta)
\]
is interpreted as an operational summary: the angle that best compresses the average contextual geometry of the aggregated sample under the chosen basis. This distinction is important under heterogeneity, because a pooled descriptive target need not coincide with the parameter favored by local predictive criteria or by participant-level generative recovery.

\subsection{Profile criterion and small-angle behavior}

Let $y_i \in \mathbb{R}^{T_i}$ denote the standardized binned latency vector for curve \(i\), and let
\[
X_i(\theta)=
\begin{bmatrix}
1 & \cos(\omega x_{i1}) & \sin(\omega x_{i1})\\
\vdots & \vdots & \vdots\\
1 & \cos(\omega x_{iT_i}) & \sin(\omega x_{iT_i})
\end{bmatrix},
\qquad \omega=\frac{\pi}{180}\theta .
\]
For fixed $\theta$, participant-level coefficients are nuisance terms estimated by ridge-stabilized least squares:
\[
\hat{\beta}_i(\theta)=
\bigl(X_i(\theta)^\top X_i(\theta)+\lambda I\bigr)^{-1}
X_i(\theta)^\top y_i,
\qquad \lambda=10^{-5}.
\]
The pooled residual criterion is
\[
RSS_\lambda(\theta)=
\sum_{i=1}^N
\left\|y_i-X_i(\theta)\hat{\beta}_i(\theta)\right\|_2^2 .
\]
We define normalized profile weights over a fixed grid of candidate angles:
\[
w(\theta\mid y)=
\frac{RSS_\lambda(\theta)^{-n/2}}
{\sum_{\theta'\in\Theta}RSS_\lambda(\theta')^{-n/2}},
\qquad n=\sum_i T_i .
\]
These weights correspond to a plug-in Gaussian profile likelihood with \(\sigma^2\) integrated under \(p(\sigma^2)\propto 1/\sigma^2\). We report the profile mode, profile-weighted mean, and highest-density intervals of \(w(\theta\mid y)\) as summaries of the task-level angular objective.

The low-angle regime has a simple analytic structure. For small \(\omega\),
\[
\cos(\omega x)=1-\frac{\omega^2 x^2}{2}+O(\omega^4), \qquad
\sin(\omega x)=\omega x+O(\omega^3).
\]
Thus, as $\theta \to 0$, the design columns approach
\[
1,\qquad 1-\frac{\omega^2 x^2}{2},\qquad \omega x.
\]
The intercept and cosine columns become nearly collinear, and the sine column contributes a weak linear component. Under sparse binning, this geometry produces weak local discrimination among low angles. This small-angle structure explains the strong competition from smooth additive baselines and the upward recovery bias observed in the calibration analyses.

\subsection{Predictive model comparison}
We compare the full angular-profile model against five baselines: (i) a cosine-only boundary model, (ii) a sine-only boundary model, (iii) a quadratic polynomial $(a+d_1x+d_2x^2)$, (iv) an exponential decay $(a+de^{-kx})$, and (v) a power law $(a+dx^p)$. These comparators serve two purposes. First, they test whether the pooled angular summary yields competitive local interpolation. Second, and more importantly, they test whether the geometric basis captures structural regularities that smooth additive models can approximate pointwise but do not represent directly. Throughout, local predictive performance is treated as a diagnostic objective distinct from the pooled profile criterion.

All models are fitted via closed-form OLS with a $10^{-5}$ ridge penalty. We measure out-of-sample performance primarily through leak-free hold-out RMSE, supplemented by approximate WAIC where computationally feasible. To ensure the findings are robust to preprocessing choices, we repeat all evaluations across the 4-, 6-, and 8-bin resolutions.

\paragraph{Why the profiled angle need not minimize hold-out RMSE.}
For each participant $i$, let
\[
\hat{\beta}_i^{\,\mathrm{train}}(\theta)
=
\arg\min_\beta
\left\|
y_i^{\mathrm{train}}-X_i^{\mathrm{train}}(\theta)\beta
\right\|_2^2.
\]
The pooled profiled estimator targets
\[
\hat{\theta}_{\mathrm{pool}}
=
\arg\min_\theta
\sum_{i=1}^N
\left\|
y_i-X_i(\theta)\hat{\beta}_i(\theta)
\right\|_2^2,
\]
or equivalently the maximizer of the profiled posterior under the Gaussian model. By contrast, temporal prediction evaluates
\[
RMSE_{\mathrm{test}}(\theta)
=
\left[
\frac{1}{\sum_i m_i}
\sum_{i=1}^N
\left\|
y_i^{\mathrm{test}}-
X_i^{\mathrm{test}}(\theta)\hat{\beta}_i^{\,\mathrm{train}}(\theta)
\right\|_2^2
\right]^{1/2},
\]
where $m_i$ is the number of held-out bins. These objectives coincide only under restrictive conditions, such as a correctly specified common angle, homogeneous dynamics, and alignment between global structural compression and short-horizon extrapolation. In heterogeneous data, the pooled criterion estimates an average structural orientation, whereas hold-out RMSE rewards local interpolation of the final bins. Boundary or reduced models may therefore optimize prediction without recovering the pooled structural summary.

\subsection{Estimator calibration and diagnostic checks}

We evaluated the pooled profiling procedure with diagnostics aimed at separating descriptive stability from latent recoverability. First, we generated synthetic interference curves at known values of $\theta$ and passed them through the full estimation pipeline to quantify recovery accuracy under controlled conditions. Second, we generated quadratic null curves to test whether concentrated posteriors over $\theta$ could arise from generic smooth temporal structure rather than from geometry specific to the interference model. Third, we permuted the temporal order within each participant curve, thereby preserving marginal latency distributions while eliminating within-curve temporal organization. Fourth, we performed bootstrap resampling over participants to assess the stability of the pooled maximum a posteriori estimate. Fifth, we examined recovery as a function of sample size to characterize how estimator behavior changes with the amount of pooled information.

Because the empirical profiles in the public tasks lie in a low-angle regime, we conducted domain-specific calibration analyses in the Gender--Science and Sexuality datasets. These analyses tested whether low generating angles could be recovered under empirically realistic coefficient and noise structures. The calibration target was practical identifiability in the region occupied by the observed data, complementing the formal full-rank argument for the basis functions.

\paragraph{Reliability and stability of derived variables.}
Reliability information was evaluated for the primary derived angular variable through
participant bootstrap resampling, sensitivity to binning resolution, D-score-like stratified
profiling, null controls, and simulation-based calibration. For conventional response-time
variables, reliability was assessed indirectly through large-sample stage-level distributional
summaries, critical-block preprocessing checks, and cross-domain replication of the same
profiling workflow. The D-score-like index was used only for descriptive stratification and
was not treated as a primary psychometric endpoint.

\subsection{Validation datasets and external comparisons}

We next examined whether the inferred geometry persisted beyond the primary pooled analysis using four external comparisons. A dataset-level summary of sources, sample sizes, retained critical blocks, latency-cleaning rules, and analytic roles is provided in Supplementary Table~S1 to make the public-domain and independent validation datasets directly comparable.

Within the Gender--Science corpus, we first stratified participants by a D-score-like congruency index, defined as the difference between mean incongruent and congruent critical-block latencies divided by the pooled critical-block standard deviation, and reprofiled $\theta$ independently within each quintile. This analysis tests whether the angular summary is merely a reformulation of a standard congruency effect-size measure or instead remains stable across substantial variation in conventional task effects.

We then applied the identical pipeline to a matched 2019 Sexuality IAT drawn from the same publicly released Project Implicit archive family. After excluding one malformed file (\texttt{iat2.0003.txt}), we processed the remaining trial-level logs using the same critical blocks, binning scheme, preprocessing rules, posterior profiling, and temporal hold-out evaluation. This comparison preserves the broad online IAT architecture, log format, and analysis pipeline while changing the construct domain. Construct content, stimulus materials, and participant composition therefore vary together, so the comparison is treated as a domain-level replication rather than as an isolated manipulation of construct content.

We next added the Age IAT as a third conventional social-category IAT. Trial-level Age IAT data were downloaded automatically from the Project Implicit OSF archive. We used the 2019 Age IAT raw archive, retained the standard critical blocks 3, 4, 6, and 7, and applied the same downstream geometric pipeline used for the other public IAT domains. Latencies were cleaned using an IAT-compatible window of 300--10{,}000 ms. Raw latency was preserved for quality-control reporting, but all downstream participant-level summaries, binned curves, posterior profiles, and geometric analyses used the cleaned latency variable. The Age IAT was included because it preserves the standard group-based target-category structure shared by conventional public IATs, thereby improving construct comparability across the public-domain analyses.

Finally, we analyzed a structurally distinct 2024 life-satisfaction IAT administered via Testable. This task differed materially from the Project Implicit corpora in platform, protocol, stimulus format, participant sample, and construct target. Unlike the public social-category IATs, the life-satisfaction task used self-referential and introspective content rather than socially shared category contrasts. From 187 raw exports, we retained only the first attempt for repeated participants, excluded cases with excessive fast responses ($>10\%$ of trials under 300 ms), and collapsed correction loops by summing latencies, yielding 180 valid participants. For comparability with the congruent--incongruent contrast used in the Project Implicit analyses, we profiled $\theta$ using only the combined critical blocks (Blocks 3 and 5). In this external dataset, the within-block coordinate was defined by actual display order within the critical block. This comparison is external rather than matched, so the life-satisfaction estimate is interpreted as a lower-precision generalization check rather than as an isolated manipulation of construct content.

\subsection*{Life-Satisfaction Implicit Association Test (IATSWB)}

The independent life-satisfaction dataset came from a 2024 Brazilian study designed to compare explicit and implicit indicators of subjective well-being under different individualism--collectivism profiles. The study was approved by the local Research Ethics Committee (CAAE 68745723.1.0000.5334). Sample size planning for the original experiment targeted two independent groups and indicated a minimum of 167 participants. Initially, 762 participants completed the self-report stage, 184 proceeded to the experimental stage, and 4 were excluded because their individualism and collectivism scores were tied, yielding a final analytic sample of 180 participants. The final sample comprised Brazilian adults aged 18--71 years (\textit{M} = 30.92, \textit{SD} = 11.38), with 71.66\% women, divided into an Individualist group ($n=68$) and a Collectivist group ($n=112$). Participants were classified as individualists when Horizontal-Individualism plus Vertical-Individualism exceeded Horizontal-Collectivism plus Vertical-Collectivism, and as collectivists when the reverse inequality held. Recruitment was conducted online through a paid Instagram advertisement directing respondents to a Google Forms survey, after which eligible participants received a link to complete the behavioral task on the Testable platform. Detailed sociodemographic characteristics of the life-satisfaction sample are reported in the Supplement.

Implicit life satisfaction was assessed with a custom Implicit Association Test for Subjective Well-Being (IATSWB), adapted from a previous IAT framework used in the context of energy-conservation self-reports. The task was designed to capture implicit self-associations with life-satisfaction content rather than deliberate self-evaluation. The instrument used four categories: \emph{true}, \emph{false}, \emph{self}, and \emph{not-self}. Stimuli were presented as short first-person sentences in Portuguese, preserving the subjective and self-referential character of life-satisfaction judgments. True and false items were used as logical anchors, whereas self and not-self classifications were based on whether life-satisfaction statements described the participant.

The task was administered online via the Testable platform and consisted of five blocks totaling 70 trials. In Block 1, participants classified 10 sentences as \emph{true} or \emph{false}. In Block 2, they classified 10 life-satisfaction statements as \emph{self} or \emph{not-self}. Block 3 was the congruent combined block, in which \emph{true} and \emph{self} shared one response key, whereas \emph{false} and \emph{not-self} shared the other response key, across 20 trials. In Block 4, participants repeated the \emph{self}/\emph{not-self} classification with reversed key mapping over 10 trials. Block 5 was the incongruent combined block, again with 20 trials, in which the pairings were reversed relative to Block 3. In the first two training blocks, feedback was provided after each response; no feedback was given in the combined test blocks. Participants were instructed to respond as quickly as possible without sacrificing accuracy, using the two response keys positioned on the keyboard, and total administration time was approximately 3 minutes.

For the external comparison, we retained 180 valid participants after preprocessing and profiled $\theta$ using only the combined critical blocks (Blocks 3 and 5), so that the analysis was aligned with the contrast between congruent and incongruent combined conditions. The within-block coordinate was the normalized display order of trials within each critical block. The IAT outcome in the original study was summarized with the conventional D-score, computed as the difference in mean response latency between the incongruent and congruent combined blocks divided by their pooled standard deviation. Following the original preprocessing description, trials with latencies exceeding $\pm 1.5$ standard deviations from the mean were excluded before statistical analysis.

\section{Results}

\subsection{Pooled profile and descriptive stability}

The primary six-bin Gender--Science analysis illustrates the full profiling workflow. The analytic sample contained \(N=141{,}329\) participants contributing binned critical-block latency trajectories. The normalized angular profile was sharply concentrated in the low-angle regime (Fig.~\ref{fig:theta_posterior}), with profile-weighted mean \(17.29^\circ\), mode \(17.25^\circ\), and 94\% HDI \([16.50^\circ,18.25^\circ]\). Posterior-predictive curves captured the broad empirical latency profile while preserving visible bin-level deviations (Fig.~\ref{fig:ppc}). The estimate remained stable across binning resolutions and bootstrap resamples, as reported in the Supplement. The estimator therefore provided stable compression of aggregated task geometry.

\begin{figure}[htbp]
\centering
\includegraphics[width=0.75\textwidth]{theta_posterior_full.png}
\caption{Profile posterior for the pooled angular summary in the six-bin Gender--Science analysis. The posterior is sharply concentrated in a low-angle regime, supporting descriptive stability of the pooled task-level profile. Shaded band: 94\% HDI.}
\label{fig:theta_posterior}
\end{figure}

\begin{figure}[htbp]
\centering
\includegraphics[width=0.75\textwidth]{ppc_full.png}
\caption{Observed and model-predicted aggregate curves at \(\theta_{\mathrm{MAP}}=17.25^\circ\). The fitted curve captures the broad latency profile and preserves visible bin-level deviations, supporting the use of \(\theta\) as a pooled task-level structural summary.}
\label{fig:ppc}
\end{figure}

The pooled profile and temporal hold-out analyses selected distinct regions of the parameter space. The pooled angular profile peaks in the low-angle interior, while temporal hold-out RMSE is minimized at the boundary \(\theta=0^\circ\) across all three binning resolutions (Fig.~\ref{fig:objectives-vs-theta}). This pattern shows that global structural compression and short-horizon forecasting capture different aspects of the latency trajectories.

\begin{figure}[htbp]
\centering
\includegraphics[width=\textwidth]{ext_objectives_vs_theta.png}
\caption{Objective-function comparison across binning resolutions. The normalized angular profile peaks in the low-angle interior, while temporal hold-out RMSE is minimized at the boundary \(\theta=0^\circ\) for 4, 6, and 8 bins. The figure shows the distinction between global structural compression and local forecasting.}
\label{fig:objectives-vs-theta}
\end{figure}

\subsection{Cross-task comparisons}
 Response-time descriptives were computed separately for congruent and incongruent critical stages (Table~\ref{tab:rt-descriptives-stage}). Because response-time distributions were right-skewed, medians and interquartile ranges are reported in the table alongside means and standard deviations. Stage-level mean response times were higher for incongruent than congruent trials in Gender--Science, Age, and the life-satisfaction IAT, but not in the Sexuality dataset. In Gender--Science, mean response time was 983.22 ms in congruent trials and 1022.69 ms in incongruent trials. In Sexuality, the mean was slightly lower for incongruent trials, decreasing from 1092.82 ms to 1073.84 ms. In Age, mean response time increased from 935.99 ms to 954.14 ms. The life-satisfaction IAT showed the largest absolute stage-level difference, increasing from 2515.39 ms in the congruent stage to 3077.98 ms in the incongruent stage.

\begin{table}[htbp]
\centering
\caption{Response-time descriptives by task and critical-stage congruency. Values are in milliseconds.}
\label{tab:rt-descriptives-stage}
\small
\begin{tabular}{@{}llrrrrr@{}}
\toprule
Task & Stage & Trials & Mean RT & Median RT & SD & IQR \\
\midrule
Gender--Science & Congruent & 8,148,717 & 983.22 & 789.00 & 665.64 & 455.00 \\
Gender--Science & Incongruent & 7,829,890 & 1022.69 & 830.00 & 671.27 & 509.00 \\
Sexuality & Congruent & 12,122,425 & 1092.82 & 869.00 & 761.39 & 524.00 \\
Sexuality & Incongruent & 11,659,198 & 1073.84 & 868.00 & 713.80 & 522.00 \\
Age & Congruent & 9,126,455 & 935.99 & 732.00 & 677.36 & 404.00 \\
Age & Incongruent & 8,780,945 & 954.14 & 751.00 & 669.20 & 457.00 \\
Life-satisfaction IAT & Congruent & 3,600 & 2515.39 & 2001.00 & 2331.73 & 1524.75 \\
Life-satisfaction IAT & Incongruent & 3,600 & 3077.98 & 2376.00 & 2863.49 & 1881.25 \\
\bottomrule
\end{tabular}
\tablenote{IAT = Implicit Association Test; SD = standard deviation; IQR = interquartile range. Response times are reported in milliseconds.}
\end{table}

D-score-like stratification provided a direct comparison with conventional effect-size variation. Among participants with a valid corrected D-score-like index, quintile means ranged from \(-0.320\) in the lowest quintile to \(0.825\) in the highest. Across these strata, the posterior over \(\theta\) remained anchored in the low-angle regime: quintile MAPs ranged only from \(17.00^\circ\) to \(17.75^\circ\) (Table~\ref{tab:theta-dlike-quintiles}). The pooled angular summary was therefore stable across large differences in conventional congruency.

\begin{table}[htbp]
\centering
\caption{Gender--Science pooled \(\theta\) by D-score-like congruency quintile.}
\label{tab:theta-dlike-quintiles}
\small
\begin{tabular}{lrrrrrrr}
\toprule
Quintile & N & Mean d-like & Theta mean & Theta SD & Theta MAP & HDI low & HDI high \\
\midrule
Q1 & 27300 & -0.684 & 16.884 & 1.007 & 17.000 & 15.000 & 18.000 \\
Q2 & 27299 & -0.242 & 17.139 & 1.044 & 17.000 & 16.000 & 19.000 \\
Q3 & 27299 & 0.069 & 17.466 & 1.113 & 17.000 & 16.000 & 20.000 \\
Q4 & 27299 & 0.374 & 17.706 & 1.171 & 18.000 & 16.000 & 20.000 \\
Q5 & 27300 & 0.788 & 17.960 & 1.249 & 18.000 & 16.000 & 20.000 \\
\bottomrule
\end{tabular}

\tablenote{MAP = maximum a posteriori estimate; HDI = highest density interval. D-score-like quintiles were defined from the corrected congruency index. Angles are reported in degrees.}
\end{table}

The matched Sexuality IAT provides a direct test of generalization across conventional public IAT domains. Processed through an identical pipeline, this matched domain yields a similarly concentrated low-angle posterior. Across 207{,}366 pooled curves in the six-bin analysis, the Sexuality posterior mean is $19.12^\circ$ (MAP $19.0^\circ$, 94\% HDI $[18.25^\circ,20.0^\circ]$; Table~\ref{tab:public-and-self-theta}; Fig.~\ref{fig:matched-public-theta}). The shift relative to Gender--Science is statistically credible (mean difference $1.82^\circ$, 94\% HDI $[0.50^\circ,3.00^\circ]$) but small in magnitude. Furthermore, the two tasks exhibit parallel behavior under different binning resolutions, showing identical downward shifts in MAP estimates as binning becomes finer.

The Age IAT provides an additional conventional social-category comparison. The raw trial-level Age dataset contained 30{,}168{,}625 rows. After restriction to critical blocks 3, 4, 6, and 7 and enforcement of the 300--10{,}000 ms latency window, 17{,}907{,}400 retained critical-block trials were available, encompassing 156{,}005 participants. The overall error rate was 0.0727. Pre-cleaning fast-trial ($<300$ ms) and long-trial ($>10{,}000$ ms) proportions were 0.0189 and 0.0012, respectively. After enforcing the latency window, the post-cleaning latency distribution had a median of 741.00 ms, a mean of 944.89 ms, a standard deviation of 673.43 ms, and an IQR of 430.00 ms.

Under the geometric pipeline, the Age IAT yielded a contextual-overlap MAP parameter of $\theta_{\text{MAP}}=18.50^\circ$, with posterior mean $18.48^\circ$ and 94\% HDI $[17.50^\circ,19.25^\circ]$. Thus, Gender--Science, Sexuality, and Age all occupy a narrow low-angle regime. This supports the interpretation that conventional public social-category IATs share a stable pooled geometry, with modest domain-specific shifts rather than qualitatively distinct angular regimes.

The independent life-satisfaction IAT provided a smaller external self-referential comparison. Using display order as the within-block coordinate, the posterior had mean \(28.49^\circ\), MAP \(17.75^\circ\), and 94\% HDI \([8.50^\circ,53.75^\circ]\) across 180 participants. The posterior was broader than in the large public corpora, consistent with the smaller sample size and different implementation, while the modal estimate remained in the low-angle range observed for Gender--Science, Sexuality, and Age.

\begin{figure}[htbp]
\centering
\includegraphics[width=0.80\textwidth]{matched_public_domain_theta_overlay.png}
\caption{Posterior comparison across public social-category IATs. Gender--Science, Sexuality, and Age all produce sharply concentrated low-angle posteriors, with modest domain-specific shifts across task domains.}
\label{fig:matched-public-theta}
\end{figure}

\begin{figure}[htbp]
\centering
\includegraphics[width=0.88\textwidth]{domain_theta_social_and_life_satisfaction.png}
\caption{Pooled angular summaries across three public social-category IATs and an independent self-referential life-satisfaction IAT. Gender--Science, Sexuality, and Age cluster in a narrow low-angle regime. The life-satisfaction IAT has a low-angle modal estimate but broader posterior uncertainty, consistent with its smaller sample size and different task implementation.}
\label{fig:domain-theta-contrast}
\end{figure}

\begin{table}[htbp]
\centering
\caption{Pooled \(\theta\) summaries across public social-category and independent life-satisfaction IATs.}
\label{tab:public-and-self-theta}
\small
\begin{tabular}{@{}lrrrrrr@{}}
\toprule
Task & Participants & Mean \(\theta\) & SD \(\theta\) & MAP \(\theta\) & HDI low & HDI high \\
\midrule
Gender--Science & 141{,}329 & 17.29 & 0.48 & 17.25 & 16.50 & 18.25 \\
Sexuality & 207{,}366 & 19.12 & 0.53 & 19.00 & 18.25 & 20.00 \\
Age & 156{,}005 & 18.48 & 0.52 & 18.50 & 17.50 & 19.25 \\
Life-satisfaction IAT & 180 & 28.49 & 13.68 & 17.75 & 8.50 & 53.75 \\
\bottomrule
\end{tabular}
\tablenote{IAT = Implicit Association Test; MAP = maximum a posteriori estimate; HDI = highest density interval; SD = standard deviation. Angles are reported in degrees.}
\end{table}

\subsection{Calibration and low-angle identifiability}

Simulation-based calibration was used to evaluate whether the profiled angle could be recovered in the empirical low-angle regime. Synthetic datasets were generated at the empirical task-level estimates and analyzed with the same preprocessing, binning, profiling, and posterior-summary pipeline used for the observed data. For the Gender--Science corpus, simulations generated at $\theta_{\mathrm{true}}=17.25^\circ$ yielded mean recovered MAPs of $39.04^\circ$ at $N=10{,}000$ and $39.62^\circ$ at $N=30{,}000$, with 94\% HDI coverage equal to 0\% in both cases. The matched Sexuality corpus produced the same qualitative pattern near $\theta_{\mathrm{true}}=19.12^\circ$.

The calibration results define the current operating range of the estimator. In the empirical low-angle regime, the pooled profile is highly stable across large samples, resampling, and task domains, while simulation-based recovery is biased upward under the present sparse-binning design. This pattern follows directly from the small-angle expansion: as $\theta \to 0$, the angular basis approaches a weakly informative polynomial-like basis, the intercept and cosine columns become nearly collinear, and the sine term contributes a weak linear component. With sparse participant-level binning, the profile has limited practical low-angle recoverability. We therefore report \(\theta\) as a task-level descriptor of contextual separation and use calibration diagnostics to define its inferential range.

\subsection{Predictive validation and objective mismatch}
On the local task of predicting held-out bins in the Gender--Science corpus, smooth additive baselines were highly competitive. The full angular-profile model and the quadratic baseline were nearly tied on temporal hold-out RMSE (21.919 vs.\ 21.924), and reduced smooth models achieved slightly lower RMSE values (Fig.~\ref{fig:model-comparison-bins6}; Table~\ref{tab:model-comparison}). The same pattern appeared in the Sexuality IAT. This result confirms that short-horizon interpolation and pooled structural profiling measure different properties of the latency data.

The matched public-domain comparison showed the same predictive pattern in Gender--Science and Sexuality. In the matched six-bin Gender--Science analysis, the interference model and quadratic baseline were nearly tied on temporal hold-out RMSE (22.067 vs.\ 22.073), while reduced smooth models achieved slightly lower RMSE values. The Sexuality IAT showed the same pattern: the interference model and quadratic baseline were again nearly tied (25.015 vs.\ 25.019), while the reduced exponential, sine-only, power-law, and cosine-only models produced marginally lower RMSE values. Thus, the public-domain analyses consistently show that local predictive performance and pooled angular profiling target different aspects of the latency data.

\begin{figure}[htbp]
\centering
\includegraphics[width=0.90\textwidth]{model_comparison_submission_bins6.png}
\caption{Temporal hold-out RMSE for the six-bin analysis. Model differences are small in absolute magnitude, with reduced smooth models achieving the lowest local forecasting error.}
\label{fig:model-comparison-bins6}
\end{figure}

\begin{table}[htbp]
\centering
\caption{Leak-free temporal hold-out comparison for the six-bin analysis (lower RMSE is better).}
\label{tab:model-comparison}
\small
\begin{tabular}{@{}lccccc@{}}
\toprule
Model & Best parameter & RMSE$_{\text{test}}$ & Mean log score$_{\text{test}}$ & \(\sigma_{\text{train}}\) & Boundary hit \\
\midrule
Interference & 124.20 & 21.919 & -21524.695 & 0.1056 & No \\
Cosine only & 180.00 & 21.899 & -7212.871 & 0.1823 & Yes \\
Sine only & 180.00 & 21.897 & -9632.215 & 0.1578 & Yes \\
Polynomial (\(2^{\text{nd}}\) order) & --- & 21.924 & -21602.217 & 0.1055 & No \\
Exponential & 10.10 & 21.896 & -9905.413 & 0.1556 & No \\
Power law & 0.05 & 21.897 & -9259.726 & 0.1609 & Yes \\
\bottomrule
\end{tabular}
\tablenote{RMSE = root mean squared error. The mean log score is reported on the held-out test bins. Boundary hit indicates whether the selected parameter reached the edge of the evaluated grid.}
\end{table}

\begin{table}[htbp]
\centering
\caption{Angular-profile hold-out robustness across numbers of bins per participant.}
\label{tab:bins-robustness}
\small
\begin{tabular}{@{}ccccc@{}}
\toprule
Bins & Selected \(\theta\) (deg) & RMSE$_{\text{test}}$ & \(\sigma_{\text{train}}\) & Boundary hit \\
\midrule
4 & 180.00 & 19.646 & 0.0000 & Yes \\
6 & 124.20 & 21.919 & 0.1056 & No \\
8 & 76.60 & 26.165 & 0.1703 & No \\
\bottomrule
\end{tabular}
\tablenote{RMSE = root mean squared error. Boundary hit indicates whether the selected parameter reached the edge of the evaluated grid. Angles are reported in degrees.}
\end{table}

\begin{table}[htbp]
\centering
\caption{Approximate WAIC comparison on the pooled full-sample fit. Higher ELPD-WAIC is better.}
\label{tab:waic-comparison}
\small
\begin{tabular}{@{}lcccc@{}}
\toprule
Model & ELPD-WAIC & \(p_{\mathrm{WAIC}}\) & SE & \(\Delta\)ELPD-WAIC \\
\midrule
Polynomial (\(2^{\text{nd}}\) order) & -810739.770 & 0.000 & 636.685 & 0.000 \\
Interference & -811813.911 & 0.012 & 632.836 & -1074.141 \\
Exponential & -918121.539 & 1.514 & 660.983 & -107381.769 \\
Power law & -927029.858 & 0.000 & 607.930 & -116290.088 \\
\bottomrule
\end{tabular}
\tablenote{WAIC = Watanabe--Akaike information criterion; ELPD = expected log predictive density; SE = standard error. Higher ELPD-WAIC values indicate better approximate out-of-sample predictive fit.}
\end{table}

\subsection{Residual and distributional fit}
Nearly identical hold-out RMSE nevertheless corresponded to different residual structures. The quadratic baseline exhibits systematic residual drift, particularly over-predicting late bins in higher-angle blocks (mean absolute residual in final bin: $0.476$ vs. $0.381$ for the interference model; Fig.~\ref{fig:residual-structure}; Table~\ref{tab:objective-mismatch}). Furthermore, the magnitude of the quadratic errors correlates more strongly with the underlying blockwise angle (\textit{r} = 0.240 vs. 0.130). Thus, the additive baseline leaves structured residual error that tracks the task's contextual geometry.

This structural deficit is even clearer in the distributional fit. Across all hold-out cells, the interference model produces more accurate full-response distributions, consistently achieving lower Wasserstein distances (mean $1.022$ vs. $1.323$), lower quantile error ($0.956$ vs. $1.283$), and lower interquartile-range error ($1.524$ vs. $2.070$; Fig.~\ref{fig:distributional-fit}; Table~\ref{tab:objective-mismatch}). Thus, additive models can match pointwise error while producing less accurate held-out distributional structure.

\begin{figure}[htbp]
\centering
\includegraphics[width=\textwidth]{residual_structure_block_holdout.png}
\caption{Held-out residual structure on block curves. Left: interference-model residual means by block and held-out bin. Middle: quadratic-model residual means. Right: quadratic minus interference residual means. The quadratic baseline leaves stronger late-bin negative drift, especially in higher-angle blocks.}
\label{fig:residual-structure}
\end{figure}

\begin{figure}[htbp]
\centering
\includegraphics[width=\textwidth]{distributional_shape_holdout.png}
\caption{Held-out distributional-fit diagnostics across the eight critical block/bin cells. Lower values are better. In every cell, the interference model yields lower Wasserstein distance, quantile error, and interquartile-range error than the quadratic baseline.}
\label{fig:distributional-fit}
\end{figure}

\begin{table}[htbp]
\centering
\caption{Model--objective-mismatch diagnostics on block-curve hold-out data. Lower is better except for the residual--\(\theta\) correlation, where values closer to zero indicate less geometry-aligned structured error.}
\label{tab:objective-mismatch}
\small
\begin{tabular}{lrr}
\toprule
Metric & Interference & Polynomial \\
\midrule
Block-curve RMSE & 57.218 & 57.235 \\
Mean abs. held-out cell residual & 0.226 & 0.252 \\
High-angle final-bin abs. residual & 0.381 & 0.476 \\
Corr(abs. residual, block $\theta$) & 0.130 & 0.240 \\
Mean Wasserstein distance & 1.022 & 1.323 \\
Mean quantile error & 0.956 & 1.283 \\
Mean IQR error & 1.524 & 2.070 \\
\bottomrule
\end{tabular}

\tablenote{RMSE = root mean squared error; IQR = interquartile range. The residual--\(\theta\) correlation quantifies the association between residual magnitude and the fitted angular summary.}
\end{table}

\subsection{Null checks}
Null analyses evaluated whether the observed posterior profile could arise from generic smoothness or disrupted temporal organization. Relative to synthetic quadratic curves, the empirical posterior was sharper and centered substantially lower, indicating that the fitted structure was not reducible to smooth additive curvature alone (Fig.~\ref{fig:null-diagnostics}; Table~\ref{tab:null-summary}). Within-curve temporal permutations produced concentrated boundary-seeking posteriors, showing that posterior sharpness must be interpreted together with temporal organization, task-domain contrasts, and calibration diagnostics.

Additional generic recovery simulations over the full angular range are reported in the Supplement. These simulations characterize estimator behavior outside the empirical low-angle range and complement the domain-specific calibration analyses reported above.

\begin{figure}[htbp]
\centering
\includegraphics[width=\textwidth]{ext_null_distributions.png}
\caption{Null posterior diagnostics versus the observed fit. The empirical posterior differs clearly from the quadratic additive null, whereas the within-curve permutation control produces concentrated boundary-seeking posteriors.}
\label{fig:null-diagnostics}
\end{figure}

\begin{table}[htbp]
\centering
\caption{Summary of posterior-shape diagnostics under the two null families.}
\label{tab:null-summary}
\small
\begin{tabular}{@{}lcccccc@{}}
\toprule
Null type & Replicates & MAP mean & MAP SD & Mean posterior SD & Mean HDI width & Mean entropy \\
\midrule
Permutation & 500 & 109.650 & 67.790 & 7.716 & 24.120 & 1.317 \\
Quadratic   & 300 & 42.250 & 13.590 & 11.591 & 41.700 & 5.200 \\
\bottomrule
\end{tabular}
\tablenote{MAP = maximum a posteriori estimate; SD = standard deviation; HDI = highest density interval. Entropy summarizes posterior dispersion under each null family.}
\end{table}

\subsection{Blockwise structure}
Blockwise profiling shows that the aggregate angle summarizes heterogeneous task phases rather than a stationary property. Profiling $\theta$ independently per block reveals systematic shifts (Fig.~\ref{fig:block-theta-summary}; Table~\ref{tab:block-theta-summary}). Block 3 yields a relatively low pooled angle ($21^\circ$), while blocks 4, 6, and 7 cluster higher ($26^\circ$--$29^\circ$). The global posterior is therefore a mixture of dynamic block-specific geometries, as contextual overlap shifts during the task.

This blockwise view further distinguishes the geometric from the additive model (Fig.~\ref{fig:block-theta-summary}). While $\theta$ moves substantially across blocks to track structural changes, blockwise quadratic coefficients merely scale in norm along a nearly constant vector in the $(d_1, d_2)$ plane. Polynomials absorb context variation through unconstrained nuisance coefficients; the geometric model maps that variation onto a rotational scale. Local predictive metrics such as RMSE are comparatively insensitive to this macroscopic organization.

\begin{figure}[htbp]
\centering
\includegraphics[width=\textwidth]{block_parameter_contrast.png}
\caption{Blockwise structure of the pooled summary. Left: block-specific profiled estimates of $\theta$ with HDI bounds. Right: block-specific quadratic coefficients in the $(d_1,d_2)$ plane. The geometric summary changes systematically across blocks, whereas the quadratic coefficients vary primarily in magnitude, supporting interpretation of the pooled angle as a structural summary over heterogeneous blockwise geometries.}
\label{fig:block-theta-summary}
\end{figure}

\begin{table}[htbp]
\centering
\caption{Block-specific pooled \(\theta\) summaries.}
\label{tab:block-theta-summary}
\small
\begin{tabular}{@{}ccccccc@{}}
\toprule
Block & Curves & Mean \(\theta\) & SD & HDI low & HDI high & MAP \\
\midrule
3 & 141258 & 20.977 & 0.828 & 19.500 & 22.500 & 21.000 \\
4 & 140027 & 29.580 & 2.388 & 25.250 & 34.000 & 29.000 \\
6 & 136487 & 26.836 & 1.807 & 23.500 & 30.250 & 26.500 \\
7 & 135700 & 26.503 & 1.751 & 23.250 & 29.750 & 26.250 \\
\bottomrule
\end{tabular}
\tablenote{SD = standard deviation; HDI = highest density interval; MAP = maximum a posteriori estimate. Angles are reported in degrees.}
\end{table}

\begin{table}[htbp]
\centering
\caption{Model comparison on block-curve hold-out data.}
\label{tab:block-model-comparison}
\small
\begin{tabular}{@{}p{5.2cm}p{4.1cm}ccc@{}}
\toprule
Model & Best parameter(s) & RMSE$_{\text{test}}$ & Mean log score$_{\text{test}}$ & Boundary hit \\
\midrule
Interference (global \(\theta\) on block curves) & 180.0 & 57.218 & -32778.299 & Yes \\
Interference (block-specific \(\theta\)) & 3: 180.0; 4: 136.8; 6: 180.0; 7: 137.5 & 57.220 & -26240.180 & Yes \\
Polynomial (block curves) & --- & 57.235 & -32781.817 & No \\
\bottomrule
\end{tabular}
\end{table}

\begin{table}[htbp]
\centering
\caption{Per-block hold-out details for the block-specific \(\theta\) selection procedure.}
\label{tab:block-theta-holdout}
\small
\begin{tabular}{@{}ccccc@{}}
\toprule
Block & Selected training \(\theta\) & RMSE$_{\text{test}}$ & Mean log score$_{\text{test}}$ & Boundary hit \\
\midrule
3 & 180.000 & 111.773 & -99475.387 & Yes \\
4 & 136.800 & 9.893 & -1438.477 & No \\
6 & 180.000 & 14.863 & -1592.495 & Yes \\
7 & 137.500 & 5.099 & -388.577 & No \\
\bottomrule
\end{tabular}
\end{table}

\subsection{Blockwise phase variation}
Block-specific phase estimates also display significant dispersion (Table~\ref{tab:block-phase-summary}). Consistent with the shifting pooled angles, this pattern indicates that aggregate contextual overlap is a macro-summary drawn across distinctly oriented block contexts.

\begin{table}[htbp]
\centering
\caption{Wrapped phase summaries for each critical block at the pooled global \(\theta_{\mathrm{MAP}}\).}
\label{tab:block-phase-summary}
\small
\begin{tabular}{@{}ccccc@{}}
\toprule
Block & Curves & Circular mean (deg) & Circular variance & Global \(\theta_{\mathrm{MAP}}\) \\
\midrule
3 & 141258 & -159.793 & 0.660 & 17.250 \\
4 & 140027 & -151.569 & 0.952 & 17.250 \\
6 & 136487 & -158.427 & 0.753 & 17.250 \\
7 & 135700 & -151.705 & 0.922 & 17.250 \\
\bottomrule
\end{tabular}
\tablenote{MAP = maximum a posteriori estimate. Circular mean and circular variance summarize wrapped phase estimates in degrees at the pooled global angular estimate.}
\end{table}

\section{Discussion}

This article contributes a diagnostic workflow for summarizing latency-trajectory geometry in large-scale behavioral tasks. Applied to three public Project Implicit IAT domains, the workflow produced a stable low-angle profile that was robust to binning choices, participant resampling, and D-score-like stratification. Applied to a smaller independent life-satisfaction IAT, it produced a compatible modal estimate with broader uncertainty. These findings support \(\theta\) as a task-level descriptor for comparing aggregated latency structure across datasets processed through the same pipeline.

The practical use case is large-scale task comparison, quality control, and exploratory characterization of latency structure. A stable pooled angular summary can reveal whether task domains occupy similar contextual-separation regimes, whether preprocessing or stimulus changes alter the shape of latency trajectories, and whether a dataset departs from the low-angle profile observed in conventional public IAT corpora. Calibration diagnostics define the estimator's operating range.

The results distinguish three properties that are often conflated in behavioral modeling: descriptive stability, local prediction, and simulation-based recoverability. On the public corpora, smooth additive baselines tie or slightly outperform the angular model on short-horizon hold-out RMSE, while the angular model shows stronger residual and distributional diagnostics. The pooled angular criterion captures a low-dimensional organization of the aggregated curves that local RMSE does not fully express.

The calibration analyses define the low-angle operating boundary. In the empirical low-angle regime, recovery is biased upward and interval coverage is limited. This behavior follows from the small-angle expansion: as \(\theta \to 0\), the angular basis approaches a weakly informative polynomial-like basis in which the intercept and cosine terms become nearly collinear and the sine term contributes a weak linear component. Under sparse binning, practical low-angle recovery is limited. A concentrated pooled profile therefore supports task-level descriptive use, with simulation-based calibration defining the range of generative interpretation.

For applied research, \(\theta\) should be reported as a pooled task-level descriptor alongside conventional IAT summaries. After identical preprocessing, researchers can compare whether IAT domains, stimulus sets, language versions, or data-collection platforms occupy similar contextual-separation regimes. The estimate should be accompanied by binning-sensitivity checks, participant bootstrap intervals, residual diagnostics, distributional-fit diagnostics, and simulation-based calibration.

The contribution is the combination of a tractable profiling estimator with explicit diagnostics for stability, objective mismatch, and calibration. The inner participant-level coefficients have closed-form least-squares solutions, the task-level angle can be profiled over a one-dimensional grid, and the resulting estimate can be stress-tested against binning changes, bootstrap resampling, additive nulls, temporal hold-out prediction, and simulation-based recovery. This combination is useful because low-dimensional summaries of behavioral time-course data can be descriptively stable even when they are not calibrated as latent-generative parameters.

The main limitations follow from the aggregation design. The analysis uses binned summaries derived from trial-level data and estimates an operational pooled target: average task-level geometry under the chosen basis. The public-domain comparisons vary construct content together with stimulus materials and participant composition, supporting domain-level comparison. The life-satisfaction comparison is smaller and differs from the Project Implicit corpora in platform, protocol, language, stimulus format, and construct target, supporting its role as an external generalization case.

Under the current estimator and data structure, \(\theta\) functions as a pooled summary of task-level contextual separation in aggregated IAT dynamics. Its stability, task sensitivity, residual diagnostics, and distributional-fit diagnostics support that role. Simulation-based calibration defines the current low-angle operating boundary.

The framework also identifies clear targets for methodological extension. Future implementations can improve low-angle conditioning through richer temporal sampling, alternative basis functions, and hierarchical regularization that models blockwise and participant-level heterogeneity directly. Empirical extensions can test whether \(\theta\) changes after controlled stimulus substitutions within the same IAT architecture, whether language- or country-specific archives show systematic shifts in contextual separation, and whether intervention or training studies alter latency geometry when conventional D-score changes are small. The present results provide a benchmark: the descriptive pattern is stable, and low-angle recovery is the key target for extending the framework toward generative inference.

\section*{Declarations}

\subsection*{Funding}
This work was financially supported by the Conselho Nacional de Desenvolvimento Cient\'ifico e Tecnol\'ogico (CNPq) under Grant No.~309589/2023-1 to A.P., and by the Coordena\c{c}\~ao de Aperfei\c{c}oamento de Pessoal de N\'ivel Superior (CAPES).

\subsection*{Competing interests}
The authors declare no competing interests.

\subsection*{Ethics approval}
The Project Implicit analyses used publicly available, anonymized archival data. The independent life-satisfaction IAT dataset came from a Brazilian study approved by the local Research Ethics Committee under CAAE 68745723.1.0000.5334.

\subsection*{Consent to participate}
Participants in the independent life-satisfaction study provided informed consent before participation.

\subsection*{Consent for publication}
Not applicable.

\subsection*{Availability of data and materials}
The public Project Implicit Gender--Science IAT raw data were obtained from the Open Science Framework at \url{https://osf.io/cfvyj/}. The public Project Implicit Sexuality IAT raw data were obtained from the Open Science Framework at \url{https://osf.io/5s9ty/}. The public Project Implicit Age IAT raw data were obtained from the Open Science Framework at \url{https://osf.io/9jvmk/}. All public IAT datasets were processed from trial-level raw logs using the critical-block preprocessing pipeline described in the Methods.

The raw Testable exports for the independent life-satisfaction IAT are not publishable as-is because they contain direct participant identifiers. The analysis database has therefore been de-identified by removing participant names, raw source filenames, exact collection timestamps, and local identifiers that could map records back to the raw exports. A de-identified life-satisfaction IAT database package has been prepared for Zenodo publication, and the dataset DOI will be inserted here after publication. To support review and reproducibility before that DOI is minted, the repository also includes documented preprocessing code, analysis scripts, and a synthetic life-satisfaction IAT dataset that reproduces the required file structure and analysis workflow without exposing participant-level raw data.

The GitHub repository and software/materials archive associated with this manuscript are published as a Zenodo software record at \url{https://doi.org/10.5281/zenodo.19711302}. Processed outputs, figure assets, reproducibility scripts, documentation, and synthetic-data materials are included in that repository archive.

\subsection*{Code availability}
Analysis code is archived in the GitHub repository's Zenodo software record at \url{https://doi.org/10.5281/zenodo.19711302}. The archive includes scripts for data ingestion, latency cleaning, participant-level binning, angular profiling, diagnostic simulation, table generation, figure export, synthetic life-satisfaction data generation, and construction of the de-identified life-satisfaction database release package.

\subsection*{Authors' contributions}
All authors contributed to study conception, analysis, and manuscript preparation. All authors approved the final version.

\appendix
\section{Profile posterior for the pooled angle}

For fixed $\theta$, let $RSS(\theta)$ denote the pooled residual sum of squares after profiling out the participant-level coefficients by least squares. Under
\[
y_i \mid \beta_i,\sigma^2,\theta \sim \mathcal{N}(X_i(\theta)\beta_i,\sigma^2 I),
\]
the Gaussian likelihood after substituting the least-squares estimates is proportional to
\[
L(\theta,\sigma^2 \mid y)\propto (\sigma^2)^{-n/2}\exp\!\left(-\frac{RSS(\theta)}{2\sigma^2}\right),
\]
where $n=\sum_i T_i$ is the total number of pooled binned observations. Using the reference prior $p(\sigma^2)\propto 1/\sigma^2$ and integrating out $\sigma^2$ gives
\[
p(\theta \mid y)\propto \int_0^\infty (\sigma^2)^{-(n/2+1)}
\exp\!\left(-\frac{RSS(\theta)}{2\sigma^2}\right)\, d\sigma^2
\propto RSS(\theta)^{-n/2}.
\]
Hence the posterior mode over $\theta$ coincides with the minimizer of the pooled profiled residual criterion.

\section{Small-angle approximation and practical identifiability}

For $\omega=\pi\theta/180$ near zero,
\[
\cos(\omega x)=1-\frac{\omega^2x^2}{2}+O(\omega^4),\qquad
\sin(\omega x)=\omega x+O(\omega^3).
\]
Therefore the interference basis $\{1,\cos(\omega x),\sin(\omega x)\}$ approaches
\[
\left\{1,\;1-\frac{\omega^2x^2}{2},\;\omega x\right\},
\]
so the intercept and cosine columns become nearly collinear while the sine column contributes a weak linear component. This produces a design that remains formally full-rank for $\omega\neq 0$ but becomes poorly conditioned as $\omega\to 0$. Under sparse binning, the profile criterion can therefore exhibit weak local discrimination among low angles, helping explain the empirical upward bias and poor interval coverage observed in the calibration analyses.

\section*{Open Practices Statement}

Materials, analysis code, processed outputs, figure assets, reproducibility scripts, and synthetic-data materials are available in the Zenodo software archive for the GitHub repository at \url{https://doi.org/10.5281/zenodo.19711302}. The public Project Implicit Gender--Science, Sexuality, and Age IAT raw datasets are linked directly in the data-availability statement. The independent life-satisfaction IAT was collected under a separate ethics protocol. Its raw Testable exports contain direct identifiers and are not publishable as-is; the analysis database has been de-identified for a separate Zenodo dataset release, with the dataset DOI to be inserted after publication. The public software archive includes documented preprocessing and analysis code, together with synthetic life-satisfaction data for code execution and review. The analyses reported in this article were not preregistered.

\bibliographystyle{apalike}
\bibliography{references}
\end{document}
